% TOP_PROBLEM: {"id": 2749, "title": "Kirby Problem 2.1", "category_id": 7, "category": {"id": 7, "name": "topology", "display_name": "Topology", "description": "Properties preserved under continuous deformations.", "slug": "topology", "order_index": 7, "created_at": "2026-07-31T15:26:25.671Z"}, "status": "open", "problem_number": "KP-2.1", "set_id": 11}
% FIRST_POSTED: 2026-10-01
% ENTRY_KIND: statement_only
% UnsolvedMath ID 2749; source catalogue code KP-2.1
% !TeX program = lualatex
% Definition and sources only; this file is not an LLM solution attempt.
\documentclass[11pt,a4paper]{article}
\usepackage[margin=25mm]{geometry}
\usepackage{fontspec}
\setmainfont{lmroman10-regular.otf}[BoldFont=lmroman10-bold.otf,
 ItalicFont=lmroman10-italic.otf,BoldItalicFont=lmroman10-bolditalic.otf]
\usepackage{amsmath,amssymb,mathtools}
\usepackage{unicode-math}
\setmathfont{latinmodern-math.otf}
\AtBeginDocument{\let\setminus\smallsetminus}
\usepackage{xurl,hyperref}
\hypersetup{colorlinks=true,urlcolor=blue}
\setcounter{secnumdepth}{0}
\setlength{\parindent}{0pt}
\setlength{\parskip}{5pt}
\setlength{\emergencystretch}{3em}
\begin{document}
\section*{Kirby Problem 2.1}\label{2749}
{\small UnsolvedMath ID 2749 · KP-2.1 · Topology · Kirby's Problems in Low-Dimensional Topology}
\subsection{Definitions and mathematical statement}
Let $S$ be a connected orientable surface of finite type and genus $g\geq3$,
with finitely many punctures and boundary components allowed.
Its mapping class group $\operatorname{Mod}(S)$ consists of isotopy classes of
orientation-preserving homeomorphisms, fixing boundary points and preserving
the puncture set; isotopies obey the same conditions.
A subgroup $G$ has finite index if there are finitely many cosets of $G$ in
$\operatorname{Mod}(S)$. Its abelianization is
\[
G^{\mathrm{ab}}=G/[G,G],
\]
where $[G,G]$ is generated by commutators $xyx^{-1}y^{-1}$.

The Ivanov conjecture asks whether $G^{\mathrm{ab}}$ is finite for every such
surface and every finite-index subgroup $G$.
These groups are finitely generated, so the question is equivalently whether
$H_1(G;\mathbb Q)=0$, or whether no such $G$ admits a nonzero homomorphism to
$\mathbb Z$. The finite abelian quotient may depend on the subgroup; there is
no proposed common upper bound for its order. Fixing punctures individually
instead changes to a finite-index subgroup and gives the same universal
finite-index formulation.
\subsection{Short English statement}
Must every finite-index subgroup of a mapping class group of genus at least three have finite abelianization?
\subsection{Sources}
\begin{itemize}
\item[S1] ULAM.AI and the UnsolvedMath Contributors, supplied problems.json, record 2749 (KP-2.1), Kirby's Problems in Low-Dimensional Topology. Catalogue identity and source transcription; CC BY 4.0.
\url{https://huggingface.co/datasets/ulamai/UnsolvedMath}.
\item[S2] K3: A New Problem List in Low-Dimensional Topology, Problem 2.1. Expanded formulation of the supplied catalogue statement.
\url{https://math.berkeley.edu/sites/default/files/surv-295-ruberman-watermarked-author-pdf.pdf}.
\end{itemize}
\end{document}
